\ExerciseSection

\begin{enumerate}
\def\labelenumi{\arabic{enumi})}
\item
  设 X 是拓扑空间，Y 是度量空间。证明在 X 的每个点处的振幅（第 IX 章，§ 2，第 3 号）都 \(\leqslant\alpha\) 的 X 到 Y 中的映射全体所成之集（其中 \(\alpha\) 是一个给定的实数 \textgreater0）在空间 \(\mathcal{F}_{u}(\mathbf{X};\mathbf{Y})\) 中是闭的。
\item
  设 X 是一个集；证明映射 \(u \to \sup_{x \in X} u(x)\) (从 \(\mathcal{B}(X; \mathbf{R})\) 到 R 中)是连续的。
\item
  设 \(\mathbf{X}\) 是度量空间，\(d\) 是 \(\mathbf{X}\) 上的度量。对每个 \(x \in \mathbf{X}\) ，设 \(d_x\) 是实值函数 \(y \to d(x, y)\) ，它在 \(\mathbf{X}\) 上连续。证明映射 \(x \to d_x\) 是 \(\mathbf{X}\) 到 \(\mathcal{C}_u(\mathbf{X}; \mathbf{R})\) 的一个子空间上的等距映射{[}该子空间赋有伪度量 \(\delta(u, v) = \sup_{x \in \mathbf{X}} |u(x) - v(x)|\) {]}。
\item
  完全正则空间 X 使得可度量化空间 \(\mathcal{C}_{u}(\mathbf{X};\mathbf{R})\) 是可数型的，当（且仅当）X 是紧且可度量化的。{[}通过考察 X 的斯通–切赫紧化（第 IX 章，§ 1，习题 7）证明 X 必是可度量化的，然后注意空间 \(\mathcal{C}_{u}(\mathbf{Z};\mathbf{R})\) 不是可数型的。{]}
\item
  设 X 是完全正则空间。
\end{enumerate}

\begin{enumerate}
\def\labelenumi{\alph{enumi})}
\tightlist
\item
  若空间 \(\mathcal{C}_{c}(\mathbf{X};\mathbf{R})\) 的每个点都有可数基本邻域系，则存在 X 的紧子集的一个递增序列 \((\mathbf{K}_{n})\) ，使得 X 的每个紧子集含于某个 \(K_{n}\) ；这时对任何可度量化空间 Y，\(\mathcal{C}_{c}(\mathbf{X};\mathbf{Y})\) 是可度量化的。
\item
  \(\mathcal{C}_{c}(\mathbf{X};\mathbf{R})\) 是可度量化且可数型的，当且仅当所有紧子空间 \(K_{n}\) 都是可度量化的（利用习题 4）。这时对每个可数型的可度量化空间 Y，\(\mathcal{C}_{c}(\mathbf{X};\mathbf{Y})\) 是可度量化且可数型的。
\end{enumerate}

\begin{enumerate}
\def\labelenumi{\arabic{enumi})}
\setcounter{enumi}{5}
\tightlist
\item
  \begin{enumerate}
  \def\labelenumii{\alph{enumii})}
  \tightlist
  \item
    设 X 是拓扑空间，Y 是度量空间，d 是 Y 上的度量，n 是整数 \textgreater0 ，S 是空间 \(X^{n} \times Y^{n} \times R_{+}^{*}\) 的一个子集。对每个点 \(z = ((x_{i}), (y_{i}), r) \in \mathbf{S}\) ，设 \(U_{z}\) 是 \(\mathcal{C}_{s}(X; Y)\) 的开子集，它由满足 \(d(f(x_i), y_i) < r\) (对 \(1 \leqslant i \leqslant n\) )的 X 到 Y 中的连续映射 f 组成。证明若 D 是 S 的稠密子集，则有 \(\bigcup_{z \in \mathbf{S}} \mathrm{U}_z = \bigcup_{z \in \mathbf{D}} \mathrm{U}_z\) 。
  \end{enumerate}
\end{enumerate}

\begin{enumerate}
\def\labelenumi{\alph{enumi})}
\setcounter{enumi}{1}
\tightlist
\item
  由 a) 推出，若 X 与 Y 的拓扑都有可数基，则 \(\mathcal{C}_{s}(\mathbf{X};\mathbf{Y})\) 的每个子空间都是林德勒夫空间（第 I 章，§ 9，习题 15），从而是仿紧的（第 IX 章，§ 4，习题 23）。
\end{enumerate}

\begin{enumerate}
\def\labelenumi{\arabic{enumi})}
\setcounter{enumi}{6}
\item
  设 X 是具有可数稠密子集的空间，Y 是每个点都有可数基本邻域系的豪斯多夫空间（相应地，可度量化空间），H 是 C(X; Y) 的一个子集。证明若 C 是 H 上的一个拓扑，它细于 X 上的逐点收敛拓扑且 H 关于它是紧的，则 H 的每个点在拓扑 C 下都有可数基本邻域系（相应地，C 是可度量化的）。（注意若 D 是 X 的任一可数稠密子集，则拓扑 C 细于 D 上的逐点收敛拓扑，且后一拓扑是豪斯多夫的。）
\item
  \begin{enumerate}
  \def\labelenumii{\alph{enumii})}
  \tightlist
  \item
    设 \(f\) 是连续映射 \((x, y) \to xy\) (从 \(\mathbf{R} \times \mathbf{R}\) 到 \(\mathbf{R}\) 中)。证明映射 \(x \to f(x, .)\) (从 \(\mathbf{R}\) 到 \(\mathcal{C}_u(\mathbf{R}; \mathbf{R})\) 中)不是连续的。
  \end{enumerate}
\end{enumerate}

\begin{enumerate}
\def\labelenumi{\alph{enumi})}
\setcounter{enumi}{1}
\tightlist
\item
  设 X、Y 是两个拓扑空间，Z 是一致空间，\(f: \mathbf{X} \times \mathbf{Y} \to \mathbf{Z}\) 是一个映射。证明若 \(f(x, .)\) 对一切 \(x \in \mathbf{X}\) 在 Y 上连续，且映射 \(x \to f(x, .)\) (从 X 到 \(\mathcal{C}_u(Y; \mathbf{Z})\) 中)连续，则 \(f\) 在 \(\mathbf{X} \times \mathbf{Y}\) 上连续。
\end{enumerate}

\begin{enumerate}
\def\labelenumi{\arabic{enumi})}
\setcounter{enumi}{8}
\tightlist
\item
  \begin{enumerate}
  \def\labelenumii{\alph{enumii})}
  \tightlist
  \item
    设 X、Y、Z 是三个拓扑空间。假设 X 与 Y 是豪斯多夫的，且这些空间中每个空间的每个点都有可数基本邻域系。设 f 是 \(X \times Y\) 到 Z 中的一个映射，使得 \(f(x, .)\) 对每个 \(x \in X\) 在 Y 上连续，且映射 \(x \to f(x, .)\) (从 X 到 \(\mathcal{C}_{c}(Y; Z)\) 中)连续。证明 f 在 \(X \times Y\) 上连续（参见 §1，第 6 号，定理 2，推论 3）。
  \end{enumerate}
\end{enumerate}

\begin{enumerate}
\def\labelenumi{\alph{enumi})}
\setcounter{enumi}{1}
\tightlist
\item
  证明当把关于 X 的假设换成 X 是局部紧的且 Z 是一致空间这一假设时，a) 的结论仍然成立（利用 § 2 习题 2 与阿斯科利定理）{[}参见习题 II b){]}。
\end{enumerate}

¶ 10) 设 X、Y 是两个拓扑空间。对 X 的每个子集 A 与 Y 的每个子集 B，设 \(T(A, B)\) 表示 \(u \in \mathcal{C}(X; Y)\) (满足 \(u(A) \subset B\) )全体所成之集。

设 \(\mathfrak{U} = (\mathbf{U}_{\alpha})\) 是 \(\mathbf{X}\) 的一个开覆盖。设 \(\mathfrak{G}_{\mathfrak{U}}\) 表示 \(\mathcal{C}(\mathbf{X};\mathbf{Y})\) 上由诸集 \(T(\mathbf{F},\mathbf{V})\) 生成的拓扑，其中 \(\mathbf{V}\) 跑遍 \(\mathbf{Y}\) 的所有开集所成之集，而 \(\mathbf{F}\) 跑遍 \(\mathbf{X}\) 的含于至少一个 \(\mathbf{U}_{\alpha}\) 中的闭子集全体所成之集。

\begin{enumerate}
\def\labelenumi{\alph{enumi})}
\item
  证明 \(\mathfrak{G}_{\mathfrak{U}}\) 细于紧开拓扑（利用第 IX 章，§ 4，第 3 号，定理 3）。若 \(\mathbf{X}\) 是正则的，则映射 \((u, x) \to u(x)\) (从 \(\mathcal{C}(\mathbf{X}; \mathbf{Y}) \times \mathbf{X}\) 到 \(\mathbf{Y}\) 中)当 \(\mathcal{C}(\mathbf{X}, \mathbf{Y})\) 赋有拓扑 \(\mathfrak{G}_{\mathfrak{U}}\) 时是连续的。
\item
  设 \(\mathfrak{C}\) 是 \(\mathcal{C}(\mathbf{X};\mathbf{Y})\) 上的一个拓扑，映射 \((u,x)\to u(x)\) (从 \(\mathcal{C}(\mathbf{X};\mathbf{Y})\times \mathbf{X}\) 到 \(\mathbf{Y}\) 中)关于它是连续的。设 \(u_0\) 是 \(\mathbf{X}\) 到 \(\mathbf{Y}\) 中的一个连续映射，\(x_0\) 是 \(\mathbf{X}\) 的一个点，\(\mathbf{V}\) 是 \(\mathbf{Y}\) 中包含 \(u_0(x_0)\) 的一个开集。证明存在一个开邻域 \(\mathbf{U}\) ( \(x_0\) 在 \(\mathbf{X}\) 中的开邻域)，使得 \(\pmb{T}(\mathrm{U},\mathrm{V})\) 是 \(u_0\) 关于 \(\mathfrak{C}\) 的一个邻域。
\item
  假设 X 是完全正则的。证明若在诸拓扑 \(\mathcal{C}\) ( \(\mathcal{C}(\mathbf{X};\mathbf{R})\) 上使 \((u,x)\to u(x)\) 连续的拓扑)中，存在一个细于所有其他拓扑的拓扑 \(\mathcal{C}_0\) ，则 \(\mathcal{C}_0\) 必是紧开拓扑。{[}设 \(T(\mathrm{U},\mathrm{V})\) 是 \(\mathcal{C}(\mathbf{X};\mathbf{R})\) 中关于 \(\mathcal{C}_0\) 的 0 的一个邻域。设 \((\mathrm{W}_{\alpha})\) 是 \(\overline{\mathbf{U}}\) 的任一开覆盖；设 \(\mathfrak{B}\) 是由诸集 \(\mathrm{W}_{\alpha}\) 与 \(\mathbb{C}\overline{\mathbf{U}}\) 组成的 X 的覆盖。利用 a) 并用反证法，证明有限多个 \(\mathrm{W}_{\alpha}\) 覆盖 \(\overline{\mathbf{U}}\) 。{]}
\item
  证明若 X 是完全正则的但不是局部紧的，则映射 \((u, x) \to u(x)\) (从 \(\mathcal{C}_{c}(\mathbf{X}; \mathbf{R}) \times \mathbf{X}\) 到 R 中)并非在每一点处连续。{[}考察一个没有紧邻域的点 \(x_{0}\) 并利用 b) 用反证法论证。{]}
\end{enumerate}

¶ 11) 设 X、Y 是两个非空拓扑空间，其乘积 \(X \times Y\) 是正规的。设 C 是集 \(\mathcal{C}(X; \mathbf{R})\) 上的一个拓扑，使得对 \(X \times Y\) 上每个连续实值函数 f，映射 \(y \to f(., y)\) (从 Y 到 \(\mathcal{C}(X; \mathbf{R})\) 中)是连续的。

\begin{enumerate}
\def\labelenumi{\alph{enumi})}
\item
  设 \(y_0\) 是 Y 的点的一个序列 \((z_n)\) (这些点都异于 \(y_0\) )的一个极限，A 是 X 的一个可数无限闭子集，其所有点都是孤立的，I 是 R 中一个有界开区间。证明关于拓扑 \(\mathfrak{G}\) ，集 \(T(\mathbf{A}, \mathbf{I})\) （习题 10）没有内点。{[}若 \(u_0 \in T(\mathbf{A}, \mathbf{I})\) ，构造一个连续映射 \(f: \mathbf{X} \times \mathbf{Y} \to \mathbf{R}\) ，使得 \(f(., y_0) = u_0\) 且 \(f(., z_n) \notin T(\mathbf{A}, \mathbf{I})\) 对所有 \(z_n \neq y_0\) 成立。{]}
\item
  由此证明，若此外映射 \((u, x) \to u(x)\) (从 \(\mathcal{C}(\mathbf{X}; \mathbf{R}) \times \mathbf{X}\) 到 R 中)关于拓扑 C 是连续的，若 X 是局部仿紧的（第 IX 章，§ 4，习题 27），且若在 X 中存在一个点的序列 \((x_{n})\) ，它收敛于一个异于所有 \(x_{n}\) 的点，则 X 必是局部紧的{[}利用习题 10 b) 与第 IX 章，§ 4，习题 25 c){]}。
\item
  由 b) 推出，若 X 是可度量化的但不是局部紧的，则 \(\mathcal{C}_{c}(\mathbf{X};\mathbf{R})\) 中有些点没有可数基本邻域系{[}利用习题 9 a){]}。利用习题 5 a) 给出一个直接证明。
\end{enumerate}

\begin{enumerate}
\def\labelenumi{\arabic{enumi})}
\setcounter{enumi}{11}
\tightlist
\item
  设 X 是拓扑空间，Y 和 Z 是两个一致空间，S 是 X 的子集的一个集，T 是 Y 的子集的一个集。
\end{enumerate}

\begin{enumerate}
\def\labelenumi{\alph{enumi})}
\item
  设 \(u_0 \in \mathcal{C}(\mathbf{X}; \mathbf{Y})\) 。证明若对每个 \(\mathbf{A} \in \mathfrak{S}\) ，\(u_0(\mathbf{A})\) 含于 \(\mathfrak{S}\) 的某个集中，则映射 \(v \to v \circ u_0\) (从 \(\mathcal{C}_{\mathfrak{T}}(\mathbf{Y}; \mathbf{Z})\) 到 \(\mathcal{C}_{\mathfrak{S}}(\mathbf{X}; \mathbf{Z})\) 中)是连续的。
\item
  设 H 是 \(\mathcal{C}_{\mathfrak{S}}(\mathrm{X};\mathrm{Y})\) 的一个子空间，\(v_{0}\in\mathcal{C}(\mathrm{Y};\mathrm{Z})\) 。证明若 \(v_{0}\) 在 \(\mathrm{H}(\mathrm{A})\) 上对所有 \(A\in S\) 一致连续，则映射 \(u\to v_{0}\circ u\) (从 H 到 \(\mathcal{C}_{\mathfrak{S}}(\mathrm{X};\mathrm{Z})\) 中)是连续的。
\item
  设 \(u_0 \in \mathcal{C}(\mathbf{X}; \mathbf{Y})\) ，\(v_0 \in \mathcal{C}(\mathbf{Y}; \mathbf{Z})\) 。假设对每个 \(\mathbf{A} \in \mathfrak{S}\) ，存在 \(\mathbf{B} \in \mathfrak{T}\) 与一个围域 \(\mathbf{V}\) ( \(\mathbf{Y}\) 的围域)满足条件：(i) \(\mathbf{V}(u_0(\mathbf{A})) \subset \mathbf{B}\) ，(ii) \(v_0\) 在 \(\mathbf{B}\) 上一致连续。
\end{enumerate}

则映射 \((u, v) \to v \circ u\) (从 \(\mathcal{C}_{\mathfrak{S}}(\mathbf{X}; \mathbf{Y}) \times \mathcal{C}_{\mathfrak{T}}(\mathbf{Y}; \mathbf{Z})\) 到 \(\mathcal{C}_{\mathfrak{S}}(\mathbf{X}; \mathbf{Z})\) 中)在 \((u_0, v_0)\) 处连续。特别地，映射 \((u, v) \to v \circ u\) (从 \(\mathcal{C}_u(\mathbf{X}; \mathbf{Y}) \times \mathcal{C}_u(\mathbf{Y}; \mathbf{Z})\) 到 \(\mathcal{C}_u(\mathbf{X}; \mathbf{Z})\) 中)在每一点 \((u_0, v_0)\) (使得 \(v_0\) 在 \(\mathbf{Y}\) 上一致连续的点)处连续。

\begin{enumerate}
\def\labelenumi{\alph{enumi})}
\setcounter{enumi}{3}
\tightlist
\item
  设 \(v_0\) 是同胚 \(x \to x^3\) (从 \(\mathbf{R}\) 到其自身上)。证明映射 \(u \to v_0 \circ u\) (从 \(\mathcal{C}_u(\mathbf{R}; \mathbf{R})\) 到其自身)并非在每一点处连续。
\end{enumerate}

¶ 13) 设 X、Y、Z 是三个豪斯多夫拓扑空间。

\begin{enumerate}
\def\labelenumi{\alph{enumi})}
\item
  证明对所有 \(u_0 \in \mathcal{C}(\mathbf{X}; \mathbf{Y})\) 与所有 \(v_0 \in \mathcal{C}(\mathbf{Y}; \mathbf{Z})\) ，映射 \(v \to v \circ u_0\) (从 \(\mathcal{C}_c(\mathbf{Y}; \mathbf{Z})\) 到 \(\mathcal{C}_c(\mathbf{X}; \mathbf{Z})\) 中)与映射 \(u \to v_0 \circ u\) (从 \(\mathcal{C}_c(\mathbf{X}; \mathbf{Y})\) 到 \(\mathcal{C}_c(\mathbf{X}; \mathbf{Z})\) 中)都是连续的。
\item
  假设 Y（相应地 Z）的每个点都有可数基本邻域系，且 X 有可数稠密子集。设 H 是 \(\mathcal{C}_{c}(\mathbf{X};\mathbf{Y})\) 的一个紧子集。证明映射 \((u,v)\to v\circ u\) (从 \(\mathbf{H}\times\mathcal{C}_{c}(\mathbf{Y};\mathbf{Z})\) 到 \(\mathcal{C}_{c}(\mathbf{X};\mathbf{Z})\) 中)是连续的。{[}先利用习题 7 与习题 9 a) 证明：若 K 是 X 的任一紧子集，则 \(\mathbf{H}(\mathbf{K})\) 是 Y 的一个紧子集。{]}
\item
  证明映射 \((u, v) \to v \circ u\) (从 \(\mathcal{C}_c(\mathbf{Q}; \mathbf{Q}) \times \mathcal{C}_c(\mathbf{Q}; \mathbf{Q})\) 到 \(\mathcal{C}_c(\mathbf{Q}; \mathbf{Q})\) 中)在任何点处都不连续。
\end{enumerate}

\begin{enumerate}
\def\labelenumi{\arabic{enumi})}
\setcounter{enumi}{13}
\tightlist
\item
  设 \(\Gamma\) 是实平面 \(\mathbf{R}^2\) 的所有同胚所成的群，赋有逐点收敛拓扑。证明映射 \((u, v) \to v \circ u\) (从 \(\Gamma \times \Gamma\) 到 \(\Gamma\) 中)不是连续的。{[}考察同胚的一个序列 \((v_n)\) ，使得 \(v_n\) 保持每一点 \((x, y)\) (满足 \(y \notin [1/(n + 1), 1/n]\) 的点)不动，且使得 \(v_{n}\) 在直线 \(y = 2/(2n + 1)\) 上的限制是平移 \((x, y) \to (x + 1, y)\) ；另一方面，考察平移的序列 \((u_{n})\)
\end{enumerate}

\[
(x, y) \rightarrow (x, y + 2 / (2 n + 1)).
\]
{]}

\begin{enumerate}
\def\labelenumi{\arabic{enumi})}
\setcounter{enumi}{14}
\tightlist
\item
  \begin{enumerate}
  \def\labelenumii{\alph{enumii})}
  \tightlist
  \item
    设 X 是一致空间，S 是 X 的子集的一个集，\(\Gamma\) 是 X 的所有同胚所成的群，\(u_{0} \in \Gamma\) 。假设对每个 \(A \in S\) 存在 \(B \in S\) 与 X 的一个围域 V，使得：\(\alpha) u_{0}^{-1}\) 在 \(\mathrm{V}(\mathrm{A})\) 上一致连续；\(\beta)\) 关系式 \(u \in \Gamma\) 与 \((u(x_{0}), u(x)) \in \mathrm{V}\) (对所有 \(x \in B\) )合起来蕴含 \(A \subset u(B)\) 。在这些条件下，证明 \(u \to u^{-1}\) (定义在 \(\Gamma\) 上)在 \(u_{0}\) 处关于 \(\Gamma\) 上的 S 收敛拓扑是连续的。
  \end{enumerate}
\end{enumerate}

\begin{enumerate}
\def\labelenumi{\alph{enumi})}
\setcounter{enumi}{1}
\tightlist
\item
  证明条件 \(\beta)\) 可以换成如下的条件：\(\beta'\) 存在一个连通集 \(\mathbf{C} \subset \mathbf{X}\) ，使得 \(\mathbf{V}(\mathbf{A}) \subset \mathbf{C}\) 且 \(\mathbf{V}(\mathbf{C})\) 含于 \(u_0(\mathbf{B})\) 的内部。{[}用反证法证明 \(\beta') \Longrightarrow \beta)\) 。{]}
\end{enumerate}

¶ 16) 设 X 是一致空间，\(\Gamma\) 是 X 的一致结构的所有自同构所成的群。

\begin{enumerate}
\def\labelenumi{\alph{enumi})}
\item
  证明 \(\Gamma\) 上的一致收敛拓扑与其群结构相容（利用习题 12 与 15）。此外，由 X 中的一致收敛一致结构在 \(\Gamma\) 上诱导的一致结构与拓扑群 \(\Gamma\) 的右一致结构一致。
\item
  若 X 是 R 的紧区间 \([0, 1]\) ，证明 a) 中定义的拓扑群 \(\Gamma\) 没有完备化{[}构造 X 的同胚的一个序列 \((u_{n})\) ，它一致收敛但使得序列 \((u_{n}^{-1})\) 不一致收敛{]}。
\item
  证明若 X 是完备豪斯多夫一致空间，则拓扑群 \(\Gamma\) 关于其双侧一致结构是完备的（第 III 章，§ 3，习题 6）。{[}注意若 \(\Phi\) 是 \(\Gamma\) 上关于这个一致结构的一个柯西滤子，则 \(\Phi(x)\) 与 \(\Phi^{-1}(x)\) 对所有 \(x \in X.\) 在 X 中收敛。{]}
\end{enumerate}

¶ 17) a) 证明若 X 是局部紧、局部连通空间，则由紧开拓扑在 X 的所有同胚所成的群 \(\Gamma\) 上诱导的拓扑与第 5 号中定义的拓扑 \(\mathcal{G}_{\beta}\) 一致{[}利用习题 15 b) 与第 5 号的命题 12{]}。

\begin{enumerate}
\def\labelenumi{\alph{enumi})}
\setcounter{enumi}{1}
\tightlist
\item
  设 X 是 R 的由点 0 与 \(2^{n} (n \in \mathbf{Z})\) 组成的局部紧子空间。若 \(\Gamma\) 是 X 的所有同胚所成的群，证明由紧开拓扑在 \(\Gamma\) 上诱导的拓扑与 \(\Gamma\) 的群结构不相容。
\end{enumerate}

\begin{enumerate}
\def\labelenumi{\arabic{enumi})}
\setcounter{enumi}{17}
\tightlist
\item
  设 X 是赋有与它的拓扑相容的一致结构的局部紧空间，\(\Gamma\) 是 X 的同胚所成的群。对 X 的每个围域 V 与 X 的每个紧子集 K，设 \(G(\mathbf{K}, \mathbf{V})\) 表示 X 的同胚的对 \((u, v)\) (满足 \((u(x), v(x)) \in \mathbf{V}\) 且 \((u^{-1}(x), v^{-1}(x)) \in \mathbf{V}\) (对所有 \(x \in K\) ))所成之集。
\end{enumerate}

\begin{enumerate}
\def\labelenumi{\alph{enumi})}
\item
  证明诸集 \(G(\mathbf{K}, \mathbf{V})\) 构成 \(\Gamma\) 上一个一致结构 U 的基本围域系，且由这个一致结构诱导的拓扑是第 5 号中定义的拓扑 \(T_{\beta}\) 。
\item
  证明若 X 关于它的一致结构是完备的，则 \(\Gamma\) 关于一致结构 U 是完备的。
\item
  证明 \(\Gamma\) 关于由拓扑 \(\mathfrak{G}_{\beta}\) 定义的双侧一致结构是完备的{[}利用习题 16 c){]}。
\item
  取 X 为 R 的由形如 \(n + 2^{-m}\) 的点 ( \(n \in Z, m\) 为整数 \(\geqslant 1\) )组成的局部紧子空间。证明在群 \(\Gamma\) 上，一致结构 U 与紧收敛一致结构都不能与由拓扑 \(C_{\beta}\) 定义的三个一致结构（左、右与双侧）中的任何一个比较。
\end{enumerate}

\begin{enumerate}
\def\labelenumi{\arabic{enumi})}
\setcounter{enumi}{18}
\tightlist
\item
  设 H 是一致空间 X 的同胚所成的一个等度连续群；赋有逐点收敛拓扑的 H 是一个拓扑群（第 5 号，命题 10 的推论）。
\end{enumerate}

\begin{enumerate}
\def\labelenumi{\alph{enumi})}
\item
  证明 H 上的左一致结构细于逐点收敛的一致结构，且若 H 是一致等度连续的，则这两个一致结构一致。
\item
  设 u 是实直线 R 的同胚，由 \(u(x) = x + 1\) (当 \(x \leqslant 0\) )、\(u(x) = x + 1/(x + 1)\) (当 \(x \geqslant 0\) )定义。设 H 是 u 在 R 的所有同胚所成的群中生成的子群。证明 H 是等度连续的，且关于逐点收敛拓扑是离散的，但 H 上的逐点收敛一致结构不是离散一致结构（注意 \(u^{n+1}(x) - u^n(x)\) 当 n 趋于 \(+\infty\) 时的极限是 0）。
\item
  取 X 为自然整数集 N 所成的离散空间，赋有度量 d，使得 \(d(m, n) = 1\) (每当 \(m \neq n\) )，并取 H 为 N 的等距所成的群，它是一致等度连续的。证明由赋有逐点收敛拓扑的 H 所得的拓扑群没有完备化{[}用与习题 16 b) 相同的方法{]}。
\item
  假设 X 是豪斯多夫且完备的，且 H 是一致等度连续的。证明 H 上关于拓扑群 \(H\) 的双侧一致结构的柯西滤子在空间 \(\mathcal{F}_s(X;X)\) 中收敛；它们的极限点所成的集 \(H'\) 是 X 的同胚所成的一个一致等度连续群，它（赋有逐点收敛拓扑）关于其双侧一致结构是完备的；且 H 在 \(H'\) 中稠密。
\end{enumerate}

\begin{enumerate}
\def\labelenumi{\arabic{enumi})}
\setcounter{enumi}{19}
\tightlist
\item
  设 X 是 \(R^{2}\) 的由诸直线 \(y = 0, y = 1/n\) ( \(n \geqslant 1\) )组成的局部紧子空间，G 是 X 的满足下列条件的所有同胚 u 所成的群：u 在每条直线 \(y = 0, y = 1/n\) 上的限制是形如 \((x, y) \to (x + a_{y}, y)\) 的平移。证明 G 满足第 5 号定理 4 的条件，但 G 上的紧收敛拓扑与逐点收敛拓扑是不同的。
\end{enumerate}

¶ 21) 设 X 是局部紧空间，T 是由满射组成的 C (X; X) 的一个子集。设 C 是 T 上的一个拓扑，它细于逐点收敛拓扑且 T 关于它是局部紧的，并考虑二元组 (T, C) 的下列性质：

\begin{enumerate}
\def\labelenumi{(\Alph{enumi})}
\tightlist
\item
  对所有 \(u \in \mathbf{T}\) 与 \(v \in \mathbf{T}\) ，有 \(u \circ v \in \mathbf{T}\) ；且对每个 \(u \in \mathbf{T}\) ，映射 \(v \to u \circ v\) 与 \(v \to v \circ u\) (从 \(\mathbf{T}\) 到其自身)关于 \(\mathfrak{C}\) 是连续的。
\end{enumerate}

\begin{enumerate}
\def\labelenumi{\alph{enumi})}
\item
  假设 X 是紧且可度量化的，二元组 (T, C) 满足 (A)，且存在 X 的同胚所成的一个群 G，它在 T 中关于拓扑 C 稠密。证明 T 中双射所成的集 H 是 T 中一个贫集的余集。{[}首先利用习题 7 注意到 T 的每个点都有一个紧可度量化邻域（关于 C）。设 \((\mathbf{V}_{n})\) 是 G 的单位元 e 在 T 中的基本邻域系，设 \(H_{n}\) 是所有这样的 \(v \in T\) 所成的集：存在 \(u \in T\) 使得 \(v \circ u \in V_{n}\) ；注意 H 包含诸集 \(H_{n}\) 的交。{]}
\item
  证明在 \(a\) 的假设下，\(\mathbf{G} \times \mathbf{X}\) 上 (映射 \(\pi: (u, x) \to u(x)\) (从 \(\mathbf{T} \times \mathbf{X}\) 到 \(\mathbf{X}\) 中)的)限制是连续的。{[}首先证明，只需证明对每个 \(x_0 \in \mathbf{X}\) ，存在 \(u_0 \in \mathbf{H}\) 使得 \(\pi\) 在 \((u_0, x_0)\) 处连续，办法是证实这蕴含 \(\pi\) 在 \((e, x_0)\) 处连续。若 \(\mathbf{V}\) 是 \(e\) 在 \(\mathbf{T}\) 中的一个紧可度量化邻域，则接着利用 \(a\) ) 与第 IX 章，§ 5，习题 23 证明：存在 \(u_0 \in \mathbf{H} \cap \mathbf{V}\) 使得 \(\pi\) 在 \((u_0, x_0)\) 处连续。{]}
\end{enumerate}

\begin{enumerate}
\def\labelenumi{\arabic{enumi})}
\setcounter{enumi}{21}
\tightlist
\item
  设 X 是紧空间，T 是 X 的所有同胚所成的群的一个子群，赋有局部紧拓扑 \(\mathcal{C}\) ，关于它习题 21 的条件 (A) 满足。设 G 是 T 的一个可数子群，\(f\) 是 X 上的一个连续实值函数。考察连续映射 \(x \to \varphi(x)\) (从 X 到 \(\mathbf{I}^{\mathrm{G}}\) 中)，其中 \(\mathbf{I} = f(\mathbf{X}) \subset \mathbf{R}\) 且 \(\varphi(x) = (f(u(x)))_{u \in \mathbf{G}}\) 。设 Y 是 X 在 \(\varphi\) 下的像；Y 是 \(\mathbf{I}^{\mathrm{G}}\) 的一个紧可度量化子空间。
\end{enumerate}

\begin{enumerate}
\def\labelenumi{\alph{enumi})}
\item
  证明所有这样的 \(v \in T\) (满足 \(\varphi \circ v = \varphi\) )所成之集是 T 的一个子群 K，它在 T 中的正规化子包含 G。设 R 是 T 上的等价关系 \(u \circ v^{-1} \in K\) ，设 \(T'\) 是商空间 T/R，\(p: T \to T/R\) 是典范映射。证明 \(T'\) 是局部紧的，且若令 \(p(u)v = p(u \circ v)\) ，则 T 右作用于 \(T'\) ，使得 \(t' \to t'v\) 在 \(T'\) 上对所有 \(v \in T\) 连续。
\item
  设 \(\mathbf{G}' = p(\mathbf{G})\) ，\(H'\) 是 \(G'\) 在 \(T'\) 中的闭包，\(\mathbf{H} = \overline{p}^{1}(\mathbf{H}')\) 。证明若 \(u \in H\) ，则关系 \(p(v) = p(v')\) 蕴含 \(p(u \circ v) = p(u \circ v')\) ，从而 H 左作用于 \(T'\) ；此外，若令 \(up(v) = p(u \circ v)\) (对 \(u \in H, v \in T\) )，则映射 \(t' \to ut'\) 在 \(T'\) 上连续。最后，若 \(u_1, u_2\) 是 H 中使 \(p(u_1) = p(u_2)\) 的元素，则有 \(u_1t' = u_2t'\) ；于是通过过渡到商，我们定义了一个映射 \((u', t') \to u't'\) (从 \(H' \times T'\) 到 \(T'\) 中)，使得 \(t' \to u't'\) 在 \(T'\) 上对所有 \(u' \in H'\) 连续。证明若 \(u' \in H'\) 且 \(v' \in H'\) ，则 \(u'v' \in H'\) ，且这样在 \(H'\) 上定义的合成律在 \(G'\) 上诱导出一个群结构（关于它 \(G'\) 同构于 GK/K）。
\item
  若 \(u \in \mathbf{H}\) ，则关系 \(\varphi(x) = \varphi(y)\) 蕴含 \(\varphi(u(x)) = \varphi(u(y))\) ，从而存在唯一的映射 \(\tilde{u}: Y \to Y\) 使得 \(\varphi \circ u = \tilde{u} \circ \varphi\) 。证明 \(\tilde{u}\) 是连续且满射的，且 \(\tilde{u}_1 = \tilde{u}_2\) 当且仅当 \(p(u_1) = p(u_2)\) ；这使我们可以写 \(\tilde{u} = \tilde{u}'\) (其中 \(u' = p(u)\) )，并定义一个单射映射 \(\psi: u' \to \tilde{u}'\) (从 \(H'\) 到 \(C_s(Y; Y)\) 中)。证明 \(\psi\) 是连续的，且 \(\psi(u'v') = \psi(u') \circ \psi(v')\) ；用 \(\psi\) 把 \(H'\) 与 \(C_s(Y; Y)\) 的一个子集等同起来，并用 \(C'\) 表示 \(H'\) 上由 \(T'\) 的拓扑诱导的拓扑，证明 \((H', C')\) 满足习题 21 的性质 (A)。
\item
  证明若 G 赋有由 \(\mathfrak{G}\) 诱导的拓扑，则映射 \((u, x) \to u(x)\) (从 \(\mathbf{G} \times \mathbf{X}\) 到 \(\mathbf{X}\) 中)是连续的。{[}证明对每个连续实值函数 \(f\) ( \(\mathbf{X}\) 上的)，映射 \((u, x) \to f(u(x))\) 在 \(\mathbf{G} \times \mathbf{X}\) 上连续，方法是利用 \(c\) ) 与习题 21；然后应用第 IX 章，§ 1，第 5 号，命题 4。{]}
\end{enumerate}

\begin{enumerate}
\def\labelenumi{\arabic{enumi})}
\setcounter{enumi}{22}
\item
  设 X 是紧空间，T 是 \(\mathcal{C}(\mathbf{X};\mathbf{X})\) 的一个在合成律 \((u,v)\to u\circ v\) 下稳定的子集。假设 T 赋有一个拓扑 \(\mathcal{C}\) ，它细于逐点收敛拓扑且 T 关于它是局部紧的；又假设对 T 的每个可数稳定子集 S，在 \(\mathbf{S}\times \mathbf{X}\) 上 (映射 \(\pi : (u,x)\to u(x)\) (从 \(\mathrm{T}\times \mathrm{X}\) 到 X 中)的)限制是连续的。证明在这些条件下，\(\pi\) 在 \(\mathrm{T}\times \mathrm{X}\) 上连续。{[}设 V 是 T 的一个紧子集（关于 \(\mathcal{C}\) ）；证明对 V 的每个可数子集 D，D 关于 \(\mathcal{C}\) 的闭包含于 D 关于一致收敛拓扑的闭包，方法是利用定理 3 的推论 1（第 4 号）。由此推出 V 在 \(\mathcal{C}_u\) (X; X) 中是相对紧的（第 II 章，§ 4，习题 6），从而是等度连续的。{]}
\item
  设 X 是局部紧空间，T 是 X 的同胚所成的一个群，C 是 T 上的一个局部紧拓扑，它细于逐点收敛拓扑且关于它习题 21 的条件 (A) 满足。证明映射 \((u, x) \to u(x)\) (从 \(T \times X\) 到 X 中)关于 C 连续。{[}把 T 的元素扩张成 X 的一点紧化 \(X'\) 的同胚；然后利用习题 22 d) 证明习题 23 的结果可以应用。{]}
\item
  设 G 是一个群，赋有拓扑 \(\mathcal{G}\) ，关于它 G 是局部紧的，且平移 \(t \to st\) 与 \(t \to ts\) 对所有 \(s \in G\) 在 G 上连续。证明 \(\mathcal{G}\) 与 G 的群结构相容（R. 埃利斯定理）。{[}把 G 与 G 的左平移所成的群等同起来，赋有逐点收敛拓扑，从而使 G 与 G 的一点紧化 X 的同胚所成的一个群（赋有逐点收敛拓扑）等同。利用习题 24 推出，在 G 上，X 中的逐点收敛拓扑与 X 中的一致收敛拓扑相同（第 4 号，定理 3，推论 1），并利用第 5 号的命题 11。{]}
\end{enumerate}

¶ 26) a) 设 X 是局部紧一致空间，G 是 X 的同胚所成的一个群，T 是 G 在 \(\mathcal{C}_{s}\) (X; X) 中的闭包。证明 T 关于逐点收敛拓扑是紧的且是一个同胚群，当且仅当 G 在 \(\mathcal{C}_{c}\) (X; X) 中是相对紧的。{[}利用习题 12 与 24 以及定理 4 的推论（第 5 号）。{]}

\begin{enumerate}
\def\labelenumi{\alph{enumi})}
\setcounter{enumi}{1}
\tightlist
\item
  设 G 是局部紧、非紧的拓扑群，X 是它的一点紧化；G 可与 X 的同胚所成的一个群等同。证明 G 不是等度连续的，尽管它在 \(\mathcal{C}_{s}(\mathbf{X};\mathbf{X})\) 中的闭包是紧的。
\end{enumerate}

\begin{enumerate}
\def\labelenumi{\arabic{enumi})}
\setcounter{enumi}{26}
\tightlist
\item
  设 X 是豪斯多夫一致空间，G 是 X 的同胚所成的一个等度连续群。
\end{enumerate}

\begin{enumerate}
\def\labelenumi{\alph{enumi})}
\item
  一点 \(x_{0} \in X\) 称为关于 G 是概周期的，如果 \(x_{0}\) 关于 G 的轨道在 X 中是相对紧的。设 Y 是这个轨道在 X 中的（紧）闭包；则 \(u(Y) \subset Y\) 对所有 \(u \in G\) 成立。设 \(\tilde{u}\) 是 u 在 Y 上的限制（作为 \(\mathcal{C}(Y; Y)\) 的一个元素）；证明 \(\tilde{u}\) 是 Y 到其自身上的一个同胚，且闭包 \(\Gamma\) (在 \(\mathcal{C}_{u}(Y; Y)\) 中，是 G 在映射 \(u \to \tilde{u}\) 下的像的闭包)是一个在 Y 上可迁的紧同胚群（第 6 号，定理 4，推论）。
\item
  假设 X 是完备的。这时 \(x_{0}\) 关于 G 是概周期的，当且仅当对 \(x_{0}\) 在 X 中的每个邻域 V，存在有限多个元素 \(u_{i} \in G\) ，使得对所有 \(u \in G\) ，至少有一个指标 i 使 \(u_{i}^{-1}(u(x_{0})) \in V\) 。
\item
  假设 \(x_{0}\) 是概周期的，且 G 赋有与它的群结构相容的拓扑。用 a) 的记号，映射 \(u \to \tilde{u}\) (从 G 到 \(\Gamma\) 中（赋有一致收敛拓扑）)连续，当且仅当对 Y 的每对不同的点 x、y，存在 e 在 G 中的一个邻域 U 与 y 的一个邻域 V，使得关系 \(u \in U\) 蕴含 \(u(x) \notin V\) {[}利用在 \(\Gamma\) 上由 \(\mathcal{C}_{u}(Y; Y)\) 与 \(\mathcal{C}_{s}(Y; Y)\) 的拓扑诱导的拓扑相同这一事实{]}。
\end{enumerate}

\begin{enumerate}
\def\labelenumi{\arabic{enumi})}
\setcounter{enumi}{27}
\tightlist
\item
  设 G 是拓扑群，X 是 G 到 C 中的有界连续映射所成的巴拿赫空间，赋有范数
\end{enumerate}

\[
\left| \left| f \right| \right| = \sup _ {s \in \mathbf {G}} | f (s) |.
\]

对每个 \(f \in \mathbf{X}\) 与每个 \(s \in \mathbf{G}\) ，设 \(U_s f\) 表示函数 \(t \to f(s^{-1}t)\) ，它属于 \(\mathbf{X}\) ，于是诸 \(U_s\) (当 \(s\) 跑遍 \(\mathbf{G}\) 时)构成一个等距所成的群 \(\mathbf{G}'\) ( \(\mathbf{X}\) 的等距群)。一个元素 \(f \in \mathbf{X}\) 称为 \(\mathbf{G}\) 上的（左）概周期函数，如果 \(f\) 是 \(\mathbf{X}\) 的关于群 \(\mathbf{G}'\) 的概周期元素。为此，必要且充分的是：对每个 \(\varepsilon > 0\) ，存在有限多个元素 \(s_i \in \mathbf{G}\) ，具有下列性质：对每个 \(s \in \mathbf{G}\) ，至少有一个指标 \(i\) 使得 \(|f(s^{-1}t) - f(s_{i}^{-1}t)| \leqslant \varepsilon\) 对所有 \(t \in \mathbf{G}\) 成立。

\begin{enumerate}
\def\labelenumi{\alph{enumi})}
\tightlist
\item
  假设 \(f\) 是概周期的。设 \(\mathbf{Y}\) 是 \(\mathbf{X}\) 中的（ \(f\) 在 \(\mathbf{G}'\) 下的轨道的）闭包，\(\Gamma\) 为在 \(\mathcal{C}_u(\mathbf{Y};\mathbf{Y})\) 中的闭包 (即诸限制 \(V_s\) (诸 \(U_s\) 在 \(\mathbf{Y} (s\in \mathbf{G})\) 上的限制)所成之集的闭包)。对每个 \(\sigma \in \Gamma\) ，设 \(f_{\sigma}\in \mathbf{Y}\) 是 \(f\) 在 \(\sigma\) 下的像。对每个 \(s\in \mathbf{G}\) ，我们有
\end{enumerate}

\[
f _ {\sigma} (t) = U _ {s} ^ {- 1} f _ {\sigma} (s ^ {- 1} t).
\]

由此推出存在一个连续函数 \(\overline{f}\) ( \(\Gamma\) 上的)，使得 \(f(s) = \overline{f}(V_s)\) {[}利用习题 27 c){]}。

\begin{enumerate}
\def\labelenumi{\alph{enumi})}
\setcounter{enumi}{1}
\tightlist
\item
  利用 a) 证明：函数 \(f \in X\) 在 G 上是概周期的，当且仅当它形如 \(g \circ \varphi\) ，其中 \(\varphi\) 是 G 到与 G 相伴的紧群 \(G^{c}\) （《集合论》，第 IV 章，§ 3，第 3 号，例 8）中的典范映射，g 是 \(G^{c}\) 到 C 中的一个连续映射。
\item
  取 G 为加法群 R。证明若 f 在 R 上是概周期的，则对每个 \(\varepsilon > 0\) 存在一个实数 T \textgreater{} 0 ，使得 R 的每个长度为 T 的区间包含一个数 s，使得
\end{enumerate}

\[
\left| f (x + s) - f (x) \right| \leqslant \varepsilon \quad \text { for   all } \quad x \in \mathbf {R}
\]

（s 称为 f 的 “ε-周期”。）{[}利用 b) 与第 V 章，§ 1 的习题 2。{]}

\begin{enumerate}
\def\labelenumi{\arabic{enumi})}
\setcounter{enumi}{28}
\tightlist
\item
  设 X 是紧空间，G 是连续作用于 X 上的一个拓扑群，使得 G 的每个轨道在 X 中稠密。证明若对 X 上每个连续实值函数 f，存在 \(x_{0} \in X\) 使得 \(s \to f(s.x_{0})\) 在 G 上是概周期的，则 G 是等度连续的（利用 X 同胚于一个方体的闭子空间这一事实）。考察当 G 交换时的逆命题。
\end{enumerate}
% semantic-fragments: legacy-input-seam
\relax{}
